% Lesson: can my cims web space run a program, not just serve files?
% Compile with:  latexmk  (settings in .latexmkrc)

\documentclass[11pt,a4paper]{article}

\usepackage{fontspec}
\usepackage[margin=2.2cm]{geometry}
\usepackage{xcolor}
\usepackage[most]{tcolorbox}
\usepackage{listings}
\usepackage{tikz}
\usetikzlibrary{positioning,arrows.meta}
\usepackage{twemojis}
\usepackage{lesson}
\usepackage{enumitem}
\usepackage[colorlinks=true,linkcolor=blue!50!black,urlcolor=blue!50!black]{hyperref}
\usepackage{newunicodechar}
\newunicodechar{→}{\ensuremath{\rightarrow}}
\newunicodechar{✗}{\ensuremath{\times}}
\newfontfamily\zhfont{Songti TC}
\newunicodechar{願}{{\zhfont 願}}
\newunicodechar{意}{{\zhfont 意}}
\newunicodechar{原}{{\zhfont 原}}
\newunicodechar{因}{{\zhfont 因}}

\setlength{\parindent}{0pt}
\setlength{\parskip}{0.6em}
\setlength{\emergencystretch}{3em}

\renewcommand{\panetitle}{Shared pane 2.2 on \name{AlfMAC5} (visible to you)}
\newcommand{\yuanyi}{\toneD{願}\toneD{意}}
\newcommand{\yuanyin}{\toneB{原}\toneA{因}}

\newtcolorbox{result}[1]{enhanced, breakable, colback=green!22, colframe=green!50!black,
  colbacktitle=green!60, coltitle=black, title={\emo{check mark button}\ #1}, fonttitle=\bfseries, boxrule=1pt, arc=2pt}

\newtcblisting{script}{listing engine=listings, listing only, enhanced,
  colback=yellow!35, colframe=orange!70!black, boxrule=0.6pt, arc=2pt,
  colbacktitle=orange!50, coltitle=black, title={\emo{page facing up}\ \name{probe.cgi}}, fonttitle=\small\bfseries,
  listing options={style=python}}

\tikzset{
  box/.style={draw=#1!70!black, fill=#1!45, very thick, rounded corners=3pt,
    minimum height=1.1cm, text width=2.3cm, align=center, font=\footnotesize},
  unknown/.style={draw=red!85!black, fill=red!12, very thick, dashed, rounded corners=3pt,
    minimum height=1.1cm, text width=2.3cm, align=center, font=\footnotesize},
  flow/.style={->, black!60},
  lbl/.style={font=\scriptsize, black!70, above, midway},
}

\title{\emo{globe with meridians}\ Can \name{cims} run a program?\\[0.3em]
  \large From a file server to one that takes a \texttt{POST}}
\author{Lesson for Alfredo Canziani, written by Claude}
\date{Sat 17 Oct 2026}

\begin{document}
\maketitle

\section*{Why we asked}

The Chinese practice app will run in the phone's browser, served from \name{cims}. When I mismatch two words, say \yuanyi\ \textit{\toneD{yuàn}\toneD{yì}} ‘willing’ with \yuanyin\ \textit{\toneB{yuán}\toneA{yīn}} ‘reason’, the app should record the mistake somewhere Claude can later read and copy into the vault. The question was whether \name{cims} can do that on its own.

\section{The starting point: a file server}

My web space on \name{cims} is the folder \name{public\_html} in my home directory there. Whatever file I put in it, the web server (\name{Apache}) hands out to anyone who asks for its address. That request, ‘give me this file’, is called a \textbf{\texttt{GET}}. That is how the dictionary is published: \name{deploy.py} copies the files there with \name{scp}, and the phone fetches them.

\begin{center}
\begin{tikzpicture}[node distance=1.1cm]
  \node[box=aqua] (phone) {phone\\browser};
  \node[box=blue, right=2.8cm of phone] (apache) {\name{Apache}\\on \name{cims}};
  \node[box=orange, right=1.6cm of apache] (file) {\name{chinese.html}\\in \name{public\_html}};
  \draw[flow] ([yshift=4pt]phone.east) -- node[lbl, align=center]{\texttt{GET}\\give me the page} ([yshift=4pt]apache.west);
  \draw[flow] ([yshift=-4pt]apache.west) -- node[lbl, below]{the page} ([yshift=-4pt]phone.east);
  \draw[flow] (apache) -- node[lbl]{reads} (file);
\end{tikzpicture}
\end{center}

This road is one-way. The web can only \emph{read} files. Only someone logged in to \name{cims}, me or Claude over \name{ssh}, can change them.

\section{What a mistake log needs: a \texttt{POST}, and something to catch it}

To save a mistake, data has to travel the other way, from the phone to \name{cims}. Browsers can send data: that request is called a \textbf{\texttt{POST}}. But a plain file cannot receive anything. Something on \name{cims} has to take the \texttt{POST} and write it down, and that something has to be a \textbf{program} the server runs when the request arrives.

\begin{center}
\begin{tikzpicture}[node distance=1.1cm]
  \node[box=aqua] (phone) {phone\\browser};
  \node[box=blue, right=2.4cm of phone] (apache) {\name{Apache}\\on \name{cims}};
  \node[unknown, right=1.4cm of apache] (prog) {a program?\\(\name{log.cgi})};
  \node[box=orange, right=1.4cm of prog] (file) {mistakes file};
  \draw[flow] (phone) -- node[lbl, align=center]{\texttt{POST}\\\yuanyi\ ✗ \yuanyin} (apache);
  \draw[flow] (apache) -- node[lbl]{runs} (prog);
  \draw[flow] (prog) -- node[lbl]{appends} (file);
\end{tikzpicture}
\end{center}

\begin{insight}{CGI}
The old, simple way for \name{Apache} to run a program is \textbf{CGI} (Common Gateway Interface). The rules are short. A file whose name ends in \texttt{.cgi} and that is allowed to execute gets \emph{run} instead of sent. Whatever the program prints is what the browser receives. It must print a header line saying what kind of content follows, then one empty line, then the content.
\end{insight}

Whether \name{Apache} on \name{cims} is allowed to run \emph{my} programs is a setting NYU chooses, and I cannot see it from outside. So we tried it.

\section{The test}

The probe is the smallest CGI program possible. It prints the header, the empty line, and the word \texttt{hello}:

\begin{script}
#!/usr/bin/env python3
print("Content-Type: text/plain")
print()
print("hello")
\end{script}

The first line tells the system which interpreter runs the file (\name{Python}~3). \texttt{text/plain} says the content is plain text. The bare \texttt{print()} is the empty line that ends the headers.

\subsection*{Step 1: look at the script}

\begin{pane}
cd ~/Work/Scratch && clear && cat probe.cgi
\end{pane}

\name{clear} wiped the pane \emph{before} \name{cat} printed the script, so the script was on screen afterwards. The commands typed before it were what got erased.

\subsection*{Step 2: upload it and make it executable}

\begin{pane}
scp probe.cgi cims:public_html/ && ssh cims 'chmod 755 public_html/probe.cgi' && echo uploaded
\end{pane}
\begin{cmdout}
probe.cgi                                       100%   80     6.8KB/s   00:00
uploaded
\end{cmdout}

\name{scp} copies the file into \name{public\_html}. \texttt{chmod 755} sets its permissions: I may read, write and run it, everyone else may read and run it. Without the ‘run’ bit, \name{Apache} would refuse to execute it. \texttt{\&\&} runs the next command only if the previous one succeeded, so \texttt{uploaded} appears only if both worked.

\subsection*{Step 3: ask for it, as a browser would}

\begin{pane}
curl -si https://cims.nyu.edu/~atcold/probe.cgi | head -15
\end{pane}
\begin{cmdout}
|!HTTP/2 200!|
|!content-type: text/plain; charset=UTF-8!|
|~date: Sat, 17 Oct 2026 16:42:47 GMT~|
|~server: Apache/2.4.62~|
|~content-security-policy: upgrade-insecure-requests~|
|~strict-transport-security: max-age=31536000; preload;~|
|~set-cookie: USERSERVERID=p1; path=/~|
|!hello!|
\end{cmdout}

\name{curl} fetches an address from the terminal: \texttt{-s} keeps it silent about progress, \texttt{-i} also shows the response headers, the lines before the content. \texttt{head -15} keeps only the first 15 lines.

\begin{result}{\name{cims} runs my programs}
The answer is \texttt{hello}, and the content type is the \texttt{text/plain} the script printed. So \name{Apache} \emph{ran} the script and sent back its output. Had it only served files, the answer would have been the script's own text, the four \name{Python} lines, or an error page. \texttt{200} is HTTP for ‘OK’.
\end{result}

\subsection*{Step 4: remove the probe}

\begin{pane}
ssh cims 'rm public_html/probe.cgi' && curl -s -o /dev/null -w '%{http_code}\n' https://cims.nyu.edu/~atcold/probe.cgi && rm probe.cgi
\end{pane}
\begin{cmdout}
|!403!|
\end{cmdout}

This deletes the script on \name{cims}, asks for its address again, and deletes the local copy. This time \name{curl} throws the content away (\texttt{-o /dev/null}) and prints only the status code (\texttt{-w}). \texttt{403} means the server refuses: the probe is no longer answering. We did not check what that address returned \emph{before} the upload, so this lesson cannot say why \name{cims} answers 403 here rather than 404 (‘not found’).

\section{What this unlocks}

The phone can now \texttt{POST} each mistake to a small program on \name{cims}, which appends it to a file. Claude then copies the new lines into the vault over \name{ssh}. Two precautions come with it:

\begin{itemize}[leftmargin=1.2em]
  \item The program answers anyone who finds its address, so the app sends a secret word with every \texttt{POST}, and the program ignores any request without it.
  \item The mistakes file lives outside \name{public\_html}, so the web can write to it only through the program, and never read it back.
\end{itemize}

\end{document}
